\textbf{Phần 1. Trắc nghiệm lựa chọn}
\setcounterpageref{ex}{0}

\begin{ex} %[3H3N20]
Để vẽ đường tròn có bán kính cho trước, ta dùng dụng cụ nào?
\choice
{\True Com-pa}
{Thước kẻ thẳng}
{Ê-ke}
{Bút chì màu}
\loigiai{
Ta dùng com-pa để vẽ đường tròn.\\
\textbf{Đáp án của câu hỏi là A}
}
\end{ex}

\begin{ex} %[3H3N20]
Để vẽ góc vuông đỉnh $O$, ta đặt đỉnh góc vuông của thước ê-ke trùng với
\choice
{Một điểm bất kì}
{Đầu của trang giấy}
{\True Đỉnh $O$}
{Tâm hình tròn}
\loigiai{
Đặt đỉnh góc vuông của ê-ke trùng với đỉnh $O$.\\
\textbf{Đáp án của câu hỏi là C}
}
\end{ex}

\begin{ex} %[3H3H20]
Để vẽ đường tròn bán kính $3\text{ cm}$, ta mở rộng hai đầu com-pa một đoạn dài
\choice
{$1{,}5\text{ cm}$}
{\True $3\text{ cm}$}
{$6\text{ cm}$}
{$9\text{ cm}$}
\loigiai{
Mở rộng com-pa bằng đúng độ dài bán kính: $3\text{ cm}$.\\
\textbf{Đáp án của câu hỏi là B}
}
\end{ex}

\begin{ex} %[3H3H20]
Khi vẽ một hình vuông có cạnh $5\text{ cm}$ trên giấy kẻ ô vuông, mỗi cạnh của hình vuông tương ứng với mấy ô vuông (nếu mỗi ô dài $1\text{ cm}$)?
\choice
{$4$ ô}
{$6$ ô}
{$10$ ô}
{\True $5$ ô}
\loigiai{
Mỗi cạnh tương ứng với 5 ô vuông.\\
\textbf{Đáp án của câu hỏi là D}
}
\end{ex}

\begin{ex} %[3H3H20]
Muốn vẽ hình chữ nhật có chiều dài $6\text{ cm}$, chiều rộng $4\text{ cm}$, ta cần vẽ mấy góc vuông?
\choice
{$2$ góc vuông}
{\True $4$ góc vuông}
{$3$ góc vuông}
{$1$ góc vuông}
\loigiai{
Hình chữ nhật có 4 góc vuông.\\
\textbf{Đáp án của câu hỏi là B}
}
\end{ex}

\begin{ex} %[3H3H20]
Vẽ hoa văn trang trí hình tròn lồng nhau: vẽ 4 đường tròn có cùng bán kính $2\text{ cm}$ có tâm lần lượt là 4 đỉnh của hình vuông cạnh $4\text{ cm}$. Bán kính của mỗi đường tròn là
\choice
{$1\text{ cm}$}
{$3\text{ cm}$}
{$4\text{ cm}$}
{\True $2\text{ cm}$}
\loigiai{
Bán kính là $2\text{ cm}$.\\
\textbf{Đáp án của câu hỏi là D}
}
\end{ex}

\begin{ex} %[3H3V20]
Bạn Lan dùng com-pa vẽ một đường tròn có đường kính $8\text{ cm}$. Khoảng cách giữa kim và đầu bút chì của com-pa khi vẽ là
\choice
{\True $4\text{ cm}$}
{$8\text{ cm}$}
{$2\text{ cm}$}
{$16\text{ cm}$}
\loigiai{
Bán kính $r = 8 : 2 = 4\text{ cm}$.\\
\textbf{Đáp án của câu hỏi là A}
}
\end{ex}

\begin{ex} %[3H3V20]
Trên lưới ô vuông cạnh $1\text{ cm}$, vẽ hình chữ nhật dài 7 ô và rộng 3 ô. Chu vi hình chữ nhật đó là
\choice
{$10\text{ cm}$}
{$21\text{ cm}$}
{\True $20\text{ cm}$}
{$14\text{ cm}$}
\loigiai{
$(7 + 3) \times 2 = 20\text{ cm}$.\\
\textbf{Đáp án của câu hỏi là C}
}
\end{ex}

\begin{ex} %[3H3V20]
Bạn Nam vẽ một bức tranh trang trí gồm 4 hình vuông nhỏ cạnh $3\text{ cm}$ ghép thành 1 hình vuông lớn. Cạnh hình vuông lớn là
\choice
{\True $6\text{ cm}$}
{$4\text{ cm}$}
{$8\text{ cm}$}
{$9\text{ cm}$}
\loigiai{
Cạnh hình vuông lớn: $3 \times 2 = 6\text{ cm}$.\\
\textbf{Đáp án của câu hỏi là A}
}
\end{ex}

\begin{ex} %[3H3V20]
Vẽ 2 đường tròn cùng có bán kính $3\text{ cm}$ và tâm cách nhau $6\text{ cm}$. Hai đường tròn này
\choice
{Cắt nhau tại 2 điểm}
{Không chạm nhau}
{\True Tiếp xúc nhau (chạm nhau tại đúng 1 điểm)}
{Trùng nhau}
\loigiai{
Khoảng cách giữa hai tâm bằng tổng hai bán kính ($3 + 3 = 6\text{ cm}$) nên hai đường tròn tiếp xúc ngoài.\\
\textbf{Đáp án của câu hỏi là C}
}
\end{ex}

\begin{ex} %[3H3V20]
Để vẽ trang trí hình cánh hoa 4 cánh đều nhau trong hình tròn, người ta dùng dụng cụ chính là
\choice
{Thước cuộn}
{\True Com-pa và thước kẻ}
{Băng dính}
{Thước dây}
\loigiai{
Dùng com-pa và thước kẻ để vẽ.\\
\textbf{Đáp án của câu hỏi là B}
}
\end{ex}

\begin{ex} %[3H3C20]
Vẽ một hình tròn tâm $O$ bán kính $4\text{ cm}$ nội tiếp trong hình vuông. Cạnh của hình vuông đó có độ dài là
\choice
{$4\text{ cm}$}
{$12\text{ cm}$}
{$16\text{ cm}$}
{\True $8\text{ cm}$}
\loigiai{
Cạnh hình vuông bằng đường kính hình tròn: $4 \times 2 = 8\text{ cm}$.\\
\textbf{Đáp án của câu hỏi là D}
}
\end{ex}

\textbf{Phần 2. Trắc nghiệm đúng sai}
\setcounterpageref{ex}{0}

\begin{ex} %[3H3H20]
Xét tính đúng sai về các bước vẽ hình:
\choiceTF
{\True Để vẽ đường tròn tâm $O$ bán kính $3\text{ cm}$, ta đặt mũi kim com-pa tại điểm $O$.}
{\True Khi vẽ góc vuông, ta kẻ 2 cạnh theo 2 cạnh góc vuông của ê-ke.}
{Độ mở của com-pa khi vẽ đường tròn chính là độ dài đường kính.}
{\True Độ mở của com-pa khi vẽ đường tròn chính là độ dài bán kính.}
\loigiai{
\begin{itemize}
    \item Đặt mũi kim tại tâm $O$: Đúng.
    \item Kẻ theo cạnh ê-ke: Đúng.
    \item Độ mở com-pa là bán kính (không phải đường kính). Ý 3 Sai, Ý 4 Đúng.
\end{itemize}
\begin{itemize}
    \item a) \textbf{Đúng}.
    \item b) \textbf{Đúng}.
    \item c) \textbf{Sai}.
    \item d) \textbf{Đúng}.
\end{itemize}
}
\end{ex}

\begin{ex} %[3H3V20]
Bạn Hùng vẽ trang trí một hình vuông lớn cạnh $8\text{ cm}$ gồm 4 hình vuông nhỏ bằng nhau. Xét tính đúng sai:
\choiceTF
{\True Mỗi hình vuông nhỏ có cạnh dài $4\text{ cm}$.}
{\True Chu vi của mỗi hình vuông nhỏ là $16\text{ cm}$.}
{Chu vi hình vuông lớn bằng tổng chu vi 4 hình vuông nhỏ.}
{\True Chu vi hình vuông lớn là $32\text{ cm}$.}
\loigiai{
\begin{itemize}
    \item Cạnh nhỏ: $8 : 2 = 4\text{ cm}$. Đúng.
    \item Chu vi nhỏ: $4 \times 4 = 16\text{ cm}$. Đúng.
    \item Chu vi lớn là $32\text{ cm} \neq 4 \times 16 = 64\text{ cm}$. Ý 3 Sai, Ý 4 Đúng.
\end{itemize}
\begin{itemize}
    \item a) \textbf{Đúng}.
    \item b) \textbf{Đúng}.
    \item c) \textbf{Sai}.
    \item d) \textbf{Đúng}.
\end{itemize}
}
\end{ex}

\textbf{Phần 3. Trả lời ngắn}
\setcounterpageref{ex}{0}

\begin{ex} %[3H3H20]
Để vẽ đường tròn bán kính $5\text{ cm}$, độ mở hai đầu kim com-pa bằng bao nhiêu xăng-ti-mét?
\shortans{5}
\loigiai{
\textbf{Đáp số:} $5$\\
$5\text{ cm}$.
}
\end{ex}

\begin{ex} %[3H3H20]
Để vẽ đường tròn đường kính $10\text{ cm}$, độ mở hai đầu kim com-pa bằng bao nhiêu xăng-ti-mét?
\shortans{5}
\loigiai{
\textbf{Đáp số:} $5$\\
$10 : 2 = 5\text{ cm}$.
}
\end{ex}

\begin{ex} %[3H3V20]
Vẽ hình chữ nhật có chiều dài $8\text{ cm}$, chiều rộng $6\text{ cm}$. Tính chu vi hình chữ nhật đó theo xăng-ti-mét.
\shortans{28}
\loigiai{
\textbf{Đáp số:} $28$\\
$(8 + 6) \times 2 = 28\text{ cm}$.
}
\end{ex}

\begin{ex} %[3H3V20]
Vẽ hình vuông có cạnh $7\text{ cm}$. Tính chu vi hình vuông đó theo xăng-ti-mét.
\shortans{28}
\loigiai{
\textbf{Đáp số:} $28$\\
$7 \times 4 = 28\text{ cm}$.
}
\end{ex}

\begin{ex} %[3H3V20]
Vẽ 3 hình vuông cạnh $2\text{ cm}$ thẳng hàng ghép lại thành một hình chữ nhật. Tính chu vi hình chữ nhật đó theo xăng-ti-mét.
\shortans{16}
\loigiai{
\textbf{Đáp số:} $16$\\
Dài: $2 \times 3 = 6\text{ cm}$, Rộng: $2\text{ cm}$. Chu vi: $(6 + 2) \times 2 = 16\text{ cm}$.
}
\end{ex}

\begin{ex} %[3H3C20]
Vẽ 4 hình tròn bằng nhau bán kính $2\text{ cm}$ xếp thẳng hàng tiếp xúc nhau. Khoảng cách giữa tâm hình tròn thứ nhất và tâm hình tròn thứ tư là bao nhiêu xăng-ti-mét?
\shortans{12}
\loigiai{
\textbf{Đáp số:} $12$\\
Khoảng cách giữa hai tâm liên tiếp là $2 + 2 = 4\text{ cm}$. Giữa 4 tâm có 3 khoảng cách: $4 \times 3 = 12\text{ cm}$.
}
\end{ex}
